\documentclass[12pt]{article}

% Idioma e codificação
\usepackage[brazil]{babel}
\usepackage[utf8]{inputenc}
\usepackage[T1]{fontenc}

% Matemática e listas
\usepackage{amsmath, amssymb}
\usepackage{enumitem}

% Layout e cores
\usepackage{geometry}
\usepackage{graphicx}
\usepackage{xcolor}

% Cores institucionais
\definecolor{maincolor}{RGB}{0, 79, 151}
\definecolor{meuazul}{RGB}{0, 79, 151}

% Cabeçalho e rodapé
\usepackage{fancyhdr}

% Blocos
\usepackage[most]{tcolorbox}
\newtcolorbox{block}[1]{%
  colback=blue!5,
  colframe=maincolor,
  coltitle=white,
  title=#1,
  fonttitle=\bfseries,
  boxrule=1.5pt,
  arc=3mm,
  left=6pt,
  right=6pt,
  top=6pt,
  bottom=6pt
}

\newtcolorbox{resposta}[1]{%
  colback=white,
  colframe=meuazul!30,
  arc=1mm,
  height=#1,
  width=\textwidth,
  boxrule=0.8pt
}

% Margens
\geometry{margin=2.5cm}

% Fancyhdr
\pagestyle{fancy}
\fancyhf{}

\renewcommand{\headrulewidth}{2pt}
\renewcommand{\footrulewidth}{2pt}

\renewcommand{\headrule}{\hbox to\headwidth{%
  \color{meuazul}\leaders\hrule height \headrulewidth\hfill}}

\renewcommand{\footrule}{\hbox to\headwidth{%
  \color{meuazul}\leaders\hrule height \footrulewidth\hfill}}

\setlength{\headheight}{52pt}
\addtolength{\topmargin}{-20pt}

\fancyhead[L]{%
  \includegraphics[height=1.2cm]{logo.png}
}

\fancyhead[R]{%
  \raggedleft
  \textcolor{meuazul}{%
    \small
    \textbf{Lista de Exercícios — Qualidade sob a Ótica do Usuário}\\
    Me.~Sergio Souza Novak\\
    \textit{Engenharia da Qualidade e Confiabilidade}
  }
}

\fancyfoot[C]{\textcolor{meuazul}{\thepage}}

% =========================
\begin{document}
% =========================

\vspace*{0.2cm}

\begin{center}
  {\Large \textbf{Avaliando a Qualidade sob a Ótica do Usuário}}\\[0.2cm]
  \textit{Instruções: Justifique suas respostas de forma clara e objetiva.
  Para questões discursivas, escreva na caixa de resposta abaixo de cada enunciado.}
\end{center}

\vspace{0.4cm}

% -------------------------------------------------------
% QUESTÃO 1
% -------------------------------------------------------
\textbf{1.} A administração de um Shopping Center lançou um app de mapa interno para ajudar os clientes. Após 1 semana, 100 usuários avaliaram o sistema usando uma Escala Likert de 5 pontos. Veja a frequência de votos tabulada abaixo:

\vspace{0.3cm}
\begin{center}
    \renewcommand{\arraystretch}{1.2}
    \begin{tabular}{l | c | c | c | c | c}
        \hline
        \textbf{Afirmação} & \textbf{DT (1)} & \textbf{D (2)} & \textbf{N (3)} & \textbf{C (4)} & \textbf{CT (5)} \\
        \hline
        \textbf{Q1} (O visual do app é muito bonito) & 0 & 5 & 10 & 45 & 40 \\
        \textbf{Q2} (O mapa carrega rapidamente) & 5 & 10 & 25 & 40 & 20 \\
        \textbf{Q3} (Consigo achar a loja que procuro) & 40 & 35 & 15 & 10 & 0 \\
        \textbf{Q4} (As rotas sugeridas fazem sentido) & 30 & 40 & 20 & 10 & 0 \\
        \textbf{Q5} (Sinto-me seguro utilizando o app) & 0 & 0 & 15 & 35 & 50 \\
        \textbf{Q6} (As informações das lojas estão atualizadas) & 5 & 10 & 30 & 35 & 20 \\
        \textbf{Q7} (Prefiro usar o app do que o Totem) & 10 & 15 & 20 & 35 & 20 \\
        \hline
    \end{tabular}
\end{center}
\vspace{0.3cm}

\noindent Analisando os dados da tabela, responda:
\begin{enumerate}[label=\alph*)]
    \item Calcule a Média matemática (de 1 a 5) apenas da afirmação \textbf{Q3}.
    \item Diante da distribuição de votos, quais são os dois grandes gargalos críticos do sistema?
    \item A diretoria apresentou um relatório com a "Média Geral do App", agrupando tudo e chegando ao valor \textbf{3.31} (Razoável/Bom). Explique o grave erro metodológico que a equipe estaria cometendo ao focar na "Média Geral" para priorizar correções.
\end{enumerate}

\begin{resposta}{7.5cm}
\end{resposta}


% -------------------------------------------------------
% QUESTÃO 2
% -------------------------------------------------------
\textbf{2.} Quando as necessidades dos usuários não são atendidas, o produto digital sofre graves consequências. Explique como o abandono e a rejeição do sistema afetam o "valor percebido" e os "custos de suporte e refatoração" da empresa.
\begin{resposta}{4.5cm}
\end{resposta}

\vspace{0.5cm}

% -------------------------------------------------------
% QUESTÃO 3
% -------------------------------------------------------
\textbf{3.} Defina o que é um Teste de Usabilidade. Durante a observação do comportamento do usuário realizando uma tarefa (como criar uma conta), quais são as métricas e reações que o avaliador deve focar em registrar?
\begin{resposta}{4.5cm}
\end{resposta}

\vspace{0.5cm}

% -------------------------------------------------------
% QUESTÃO 4
% -------------------------------------------------------
\textbf{4.} O SUS (System Usability Scale) é amplamente utilizado pela indústria para medir a usabilidade. Descreva como funciona a sua estrutura e explique por que é uma boa prática adotar a "alternância de polaridade" (perguntas positivas intercaladas com negativas) em questionários de avaliação.
\begin{resposta}{5cm}
\end{resposta}

\vspace{0.5cm}

% -------------------------------------------------------
% QUESTÃO 5
% -------------------------------------------------------
\textbf{5.} Após calcular a pontuação final do SUS de um sistema recém-lançado, o resultado numérico obtido pela equipe foi "35". De acordo com a escala de interpretação do SUS, como você deve classificar qualitativamente esse resultado?
\begin{resposta}{3.5cm}
\end{resposta}

\vspace{0.5cm}

% -------------------------------------------------------
% QUESTÃO 6
% -------------------------------------------------------
\textbf{6.} O Net Promoter Score (NPS) foca na lealdade do usuário. Explique a lógica de agrupar os respondentes em "Promotores", "Neutros" e "Detratores", e escreva a fórmula utilizada para calcular o score final do NPS.
\begin{resposta}{4.5cm}
\end{resposta}

\vspace{0.5cm}

% -------------------------------------------------------
% QUESTÃO 7
% -------------------------------------------------------
\textbf{7.} Os Mapas de Calor (Heatmaps) são ferramentas analíticas extremamente visuais. Descreva quais tipos de interações do usuário eles rastreiam e o que representam as cores quentes (vermelho) e as cores frias (azul) no mapa.
\begin{resposta}{4.5cm}
\end{resposta}

\vspace{0.5cm}

% -------------------------------------------------------
% QUESTÃO 8
% -------------------------------------------------------
\textbf{8.} O padrão de leitura em "F" (F-Pattern) é um comportamento clássico documentado por mapas de calor. Descreva como o olho humano escaneia a tela nesse padrão e explique o fenômeno da "Cegueira de Banner" no lado direito da tela.
\begin{resposta}{5cm}
\end{resposta}

\vspace{0.5cm}

% -------------------------------------------------------
% QUESTÃO 9
% -------------------------------------------------------
\textbf{9.} A Escala Likert de 5 pontos é considerada um padrão para medir a concordância em pesquisas qualitativas. Explique o que é o "Viés de Tendência Central" e como ele pode mascarar a opinião real do respondente.
\begin{resposta}{4cm}
\end{resposta}

\vspace{0.5cm}

% -------------------------------------------------------
% QUESTÃO 10
% -------------------------------------------------------
\textbf{10.} Por que analisar estatisticamente a "Média Geral" de um questionário muitas vezes não reflete os verdadeiros gargalos de usabilidade de um sistema? Justifique por que, na Engenharia de Software, a análise das respostas a cada item individualmente é o que deve guiar a refatoração do código.
\begin{resposta}{5.5cm}
\end{resposta}

\vspace{0.5cm}

% -------------------------------------------------------
% QUESTÃO 11 (Estudo de Caso)
% -------------------------------------------------------

\vspace{0.8cm}

\begin{center}
  \textit{``Boa atividade! Att. Prof. Me. Sergio Souza Novak.''}
\end{center}

% =========================
\end{document}
% =========================
